% !TEX program = xelatex
% =====================================================================
% BÁO CÁO KIỂM THỬ PHẦN MÔ HÌNH AI - HỆ THỐNG eKYC / ID VERIFICATION
% Nguồn số liệu: ai/tests, ai/evaluation, ai/reports
% Biên dịch bằng XeLaTeX hoặc LuaLaTeX (KHÔNG dùng pdfLaTeX):
%   xelatex bao_cao_kiem_thu_ai.tex   (chạy 2 lần để cập nhật tham chiếu)
% =====================================================================
\documentclass[12pt,a4paper]{article}

% ---------- Font & ngôn ngữ ----------
\usepackage{fontspec}
\usepackage[main=vietnamese,provide=*]{babel}

% Dùng Times New Roman / Arial / Consolas nếu có; nếu không (ví dụ Overleaf)
% thì tự động chuyển sang font tương đương.
\ExplSyntaxOn
\NewDocumentCommand{\FontOrElse}{mmm}{% kiểm tra font bằng fontspec: {tên font}{nếu có}{nếu không}
  \fontspec_font_if_exist:nTF{#1}{#2}{#3}}
\ExplSyntaxOff

\FontOrElse{Times New Roman}{\setmainfont{Times New Roman}}{%
  \FontOrElse{TeX Gyre Termes}{\setmainfont{TeX Gyre Termes}}{%
  \FontOrElse{Liberation Serif}{\setmainfont{Liberation Serif}}{}}}
\FontOrElse{Arial}{\setsansfont{Arial}}{%
  \FontOrElse{TeX Gyre Heros}{\setsansfont{TeX Gyre Heros}}{%
  \FontOrElse{Liberation Sans}{\setsansfont{Liberation Sans}}{}}}
\FontOrElse{Consolas}{\setmonofont{Consolas}}{%
  \FontOrElse{DejaVu Sans Mono}{\setmonofont{DejaVu Sans Mono}[Scale=0.9]}{%
  \FontOrElse{Liberation Mono}{\setmonofont{Liberation Mono}[Scale=0.9]}{}}}

% ---------- Bố cục ----------
\usepackage[a4paper,left=3cm,right=2cm,top=2.5cm,bottom=2.5cm]{geometry}
\usepackage{setspace}
\setstretch{1.3}
\setlength{\parindent}{0pt}
\setlength{\parskip}{0.5em}
\setlength{\emergencystretch}{2em}

% ---------- Toán, bảng, danh sách ----------
\usepackage{amsmath,amssymb}
\usepackage{siunitx}
\usepackage{booktabs}
\usepackage{longtable}
\usepackage{array}
\usepackage{graphicx}
\usepackage{float}
\usepackage{enumitem}
\usepackage{xcolor}
\usepackage{caption}
\captionsetup{font=small,labelfont=bf,labelsep=period}
\renewcommand{\arraystretch}{1.15}

% ---------- Khung mã lệnh ----------
\usepackage{listings}
\definecolor{codebg}{gray}{0.96}
\lstdefinelanguage{PS}{%
  morekeywords={cd,python},
  morecomment=[l]{\#},
  morestring=[b]"}
\lstset{
  language=PS,
  basicstyle=\ttfamily\footnotesize,
  breaklines=true,
  backgroundcolor=\color{codebg},
  frame=single,
  rulecolor=\color{gray!50},
  xleftmargin=2mm,
  xrightmargin=2mm,
  keywordstyle=\color{blue!70!black},
  commentstyle=\color{gray},
  stringstyle=\color{green!40!black},
  showstringspaces=false,
}

% ---------- Liên kết ----------
\usepackage{hyperref}
\hypersetup{hidelinks}

% ---------- Tiêu đề ----------
\title{\textbf{BÁO CÁO KIỂM THỬ PHẦN MÔ HÌNH AI\\[1mm]
TRONG HỆ THỐNG XÁC THỰC DANH TÍNH ĐIỆN TỬ (eKYC)}\\[4mm]
\large Mã dự án: IDVerification / AI eKYC Service}

\author{
  \textbf{Sinh viên:} \underline{\hspace{5cm}}\\[2mm]
  \textbf{Mã sinh viên:} \underline{\hspace{4cm}}\\[2mm]
  \textbf{Khoa:} \underline{\hspace{5cm}}\\[2mm]
  \textbf{Năm học:} \underline{\hspace{4cm}}
}
\date{\today}

\begin{document}
\maketitle

\begin{abstract}
Báo cáo này tổng hợp kết quả chạy kiểm thử đơn vị (unit test) và các đánh giá định lượng đã thực hiện
trong quá trình xây dựng hệ thống xác thực danh tính điện tử (eKYC) cho căn cước công dân (CCCD) và
chứng minh nhân dân (CMND) Việt Nam. Phạm vi gồm bảy nhóm chức năng của dịch vụ AI: nhận dạng ký tự
quang học (OCR), đối chiếu mã QR, phát hiện và so khớp khuôn mặt, passive liveness, cổng chất lượng ảnh
(liveness gate), active liveness và phát hiện giấy tờ bị chỉnh sửa (tamper detection).

Toàn bộ \textbf{44/44 test} đều đạt (0 thất bại, 0 bị bỏ qua), thời gian thực thi lần gần nhất khoảng
\num{10.64} giây. Bên cạnh đó, hệ thống được đánh giá định lượng trên ba bộ dữ liệu công khai là
CASIA-FASD (passive liveness), LFW (so khớp khuôn mặt) và FantasyID (phát hiện giả mạo giấy tờ), đồng thời
đo độ trễ tuần tự của pipeline. Kết quả cho thấy passive liveness và so khớp khuôn mặt đạt hiệu năng tốt
trên dữ liệu công khai, trong khi mô-đun TamperNet chưa đạt yêu cầu trên FantasyID và cần được cải thiện.
\end{abstract}

% =====================================================================
\section{Tổng quan kết quả kiểm thử}
% =====================================================================
Bộ kiểm thử gồm \textbf{44 test}, kết quả \textbf{44 passed, 0 failed, 0 skipped} (tỷ lệ thành công 100\%).
Thời gian chạy lần gần nhất là \num{10.64} giây (lần đầu ghi nhận: \num{11.97} giây).

\begin{table}[H]
  \centering
  \caption{Tổng hợp kết quả kiểm thử đơn vị theo nhóm chức năng}
  \label{tab:summary}
  \small
  \begin{tabular}{@{}p{3.6cm}p{5.2cm}cccc@{}}
    \toprule
    \textbf{Nhóm} & \textbf{Tệp kiểm thử} & \textbf{Số test} & \textbf{PASS} & \textbf{FAIL} & \textbf{SKIP}\\
    \midrule
    OCR & \path{test_ocr_service.py} & 4 & 4 & 0 & 0\\
    QR cross-check & \path{test_qr_service.py} & 8 & 8 & 0 & 0\\
    Face detection / matching & \path{test_face_service.py} & 6 & 6 & 0 & 0\\
    Passive Liveness (MiniFASNetV2) & \path{test_liveness.py} & 4 & 4 & 0 & 0\\
    Liveness Gate / Decision & \path{test_liveness_gate.py} & 6 & 6 & 0 & 0\\
    Active Liveness & \path{test_active_liveness.py} & 4 & 4 & 0 & 0\\
    Tamper Detection & \path{test_tamper.py} & 12 & 12 & 0 & 0\\
    \midrule
    \textbf{Tổng} & & \textbf{44} & \textbf{44} & \textbf{0} & \textbf{0}\\
    \bottomrule
  \end{tabular}
\end{table}

Các test có thời gian chạy lâu nhất:
\begin{itemize}
  \item Regression khoảng giá trị đầu vào của MiniFASNetV2: \num{2.36} giây.
  \item Hợp đồng (contract) của tamper engine khi chưa nạp model: \num{1.81} giây.
  \item Bộ phát hiện khuôn mặt UltraFace chạy thực tế: \num{1.13} giây.
  \item TamperNet forward / loss: \num{0.16} giây.
\end{itemize}

% =====================================================================
\section{Phương pháp và môi trường kiểm thử}
% =====================================================================
\subsection{Môi trường thực nghiệm}
\begin{table}[H]
  \centering
  \caption{Môi trường chạy kiểm thử}
  \label{tab:env}
  \small
  \begin{tabular}{@{}ll@{}}
    \toprule
    \textbf{Thành phần} & \textbf{Giá trị}\\
    \midrule
    Hệ điều hành & Windows 10, build 26200\\
    Python & 3.11.9\\
    pytest & 9.1.1\\
    PyTorch & 2.13.0 + cu126 (có hỗ trợ CUDA)\\
    GPU & NVIDIA GeForce RTX 4060 Laptop (8 GB VRAM)\\
    ONNX Runtime & 1.24.3 (CPUExecutionProvider)\\
    OpenCV & 4.10.0\\
    NumPy & 2.3.5\\
    MediaPipe & 1.0.1\\
    Thời gian chạy toàn bộ & 44 test trong khoảng \num{10.64} giây\\
    \bottomrule
  \end{tabular}
\end{table}

\subsection{Phạm vi và tiêu chí đánh giá}
\begin{itemize}
  \item Kiểm thử đơn vị tập trung vào logic quyết định, tiền xử lý, cơ chế dự phòng (fallback) và ngưỡng.
  \item Kiểm thử tích hợp nội bộ giữa các service trong dịch vụ AI được thực hiện thông qua đối tượng giả (\texttt{Mock}).
  \item Bộ kiểm thử không yêu cầu kết nối Internet hay cơ sở dữ liệu.
  \item Một số test chạy model thật; các test còn lại chỉ kiểm tra logic thuần.
  \item Một số mô-đun chưa có test tự động riêng và được đánh giá thông qua benchmark trên dữ liệu công khai (Mục~\ref{sec:quant}).
\end{itemize}

\subsection{Các lệnh chạy kiểm thử}
Chạy toàn bộ bộ kiểm thử (PowerShell):
\begin{lstlisting}
cd ai
$env:PYTHONIOENCODING="utf-8"
python -m pytest tests -v --tb=short
\end{lstlisting}

Chạy riêng từng nhóm:
\begin{lstlisting}
cd ai
$env:PYTHONIOENCODING="utf-8"

# OCR
python -m pytest tests/test_ocr_service.py -v
# QR cross-check
python -m pytest tests/test_qr_service.py -v
# Face detection & matching
python -m pytest tests/test_face_service.py -v
# Passive liveness
python -m pytest tests/test_liveness.py -v
# Liveness gate & decision logic
python -m pytest tests/test_liveness_gate.py -v
# Active liveness
python -m pytest tests/test_active_liveness.py -v
# Tamper detection
python -m pytest tests/test_tamper.py -v
\end{lstlisting}

% =====================================================================
\section{Chi tiết kiểm thử theo nhóm chức năng}
% =====================================================================
\subsection{OCR}
\textit{Tệp kiểm thử:} \path{tests/test_ocr_service.py}

\textbf{Mục tiêu.} Kiểm tra logic trích xuất số giấy tờ từ văn bản OCR, gồm ưu tiên số CCCD 12 chữ số và
dự phòng sang CMND 9 chữ số.

\begin{table}[H]
  \centering
  \caption{Kịch bản kiểm thử nhóm OCR}
  \label{tab:ocr-results}
  \small
  \begin{tabular}{@{}p{4.2cm}p{5.6cm}p{4.2cm}@{}}
    \toprule
    \textbf{Kịch bản} & \textbf{Đầu vào} & \textbf{Kết quả mong đợi}\\
    \midrule
    Có số CCCD 12 chữ số & Văn bản OCR chứa ``Số: 012345678901'' & Trả về \texttt{012345678901}\\
    Có cả CCCD 12 và CMND 9 chữ số & ``CMND 123456789 / CCCD 012345678901'' & Ưu tiên trả về số CCCD\\
    Chỉ có CMND 9 chữ số & ``CMND số 123456789'' & Trả về số CMND\\
    Không có số giấy tờ & Văn bản không chứa số hợp lệ & Trả về \texttt{null}\\
    \bottomrule
  \end{tabular}
\end{table}

\subsection{QR cross-check}
\textit{Tệp kiểm thử:} \path{tests/test_qr_service.py}

\textbf{Mục tiêu.} Kiểm tra giải mã và phân tích (parse) payload QR, đối chiếu QR với kết quả OCR, xử lý các
trường hợp thiếu trường legacy, thiếu QR, thiếu OCR và giải mã ảnh mẫu (fixture).

\textbf{Kết quả.}
\begin{itemize}
  \item Giải mã thành công ảnh mẫu \path{tests/fixtures/cccd_back_qr.png}, thu được số giấy tờ có dạng
        \texttt{068096******} (đã che các chữ số cuối).
  \item Phân tích payload đúng cho cả hai trường hợp có và không có số CMND cũ (legacy).
  \item Đối chiếu OCR $\leftrightarrow$ QR trả về một trong bốn trạng thái: \texttt{MATCH}, \texttt{PARTIAL},
        \texttt{MISMATCH}, \texttt{NOT\_FOUND}.
  \item So sánh không phân biệt chữ hoa/thường, bỏ qua dấu thanh và ký tự phân tách.
\end{itemize}

\subsection{Face detection và matching}
\textit{Tệp kiểm thử:} \path{tests/test_face_service.py}

\textbf{Mục tiêu.} Kiểm tra kích thước đầu vào của UltraFace, tiền xử lý, hậu xử lý (NMS) và cơ chế dự phòng
khi trích xuất embedding.

\begin{itemize}
  \item Kích thước đầu vào của UltraFace phải là \texttt{(640, 480)}.
  \item Tiền xử lý tạo tensor \texttt{(1, 3, 480, 640)} kiểu float32.
  \item Hậu xử lý áp dụng NMS và loại bỏ các box có điểm thấp; giữ lại các box tách biệt.
  \item Khi không có box từ landmark, embedding dùng box do detector phát hiện.
\end{itemize}

\subsection{Passive Liveness}
\textit{Tệp kiểm thử:} \path{tests/test_liveness.py}

\textbf{Mục tiêu.} Kiểm tra hình học của bước crop khuôn mặt và một regression quan trọng: đầu vào của
MiniFASNetV2 phải nằm trong khoảng $[0, 255]$, không được chuẩn hóa về $[0, 1]$.

\begin{itemize}
  \item Box được mở rộng 2,7 lần vẫn nằm trong ảnh; ảnh crop có kích thước $80\times 80$.
  \item Hàm dự đoán trả về ba xác suất cho ba lớp: \textit{real}, \textit{print}, \textit{replay}, với tổng bằng 1.
  \item Trạng thái \texttt{is\_live} được suy ra từ lớp \textit{real}.
\end{itemize}

\subsection{Liveness Gate và quyết định trạng thái}
\textit{Tệp kiểm thử:} \path{tests/test_liveness_gate.py}

\textbf{Mục tiêu.} Kiểm tra cổng chất lượng ảnh (blur gate) và hành vi \textit{retry}/\textit{reject} của
\texttt{VerifyService}.

\begin{itemize}
  \item Ảnh quá mờ: không gọi MiniFASNet, trả về \texttt{LOW\_\allowbreak QUALITY} và yêu cầu chụp lại (\texttt{retake\_\allowbreak required}).
  \item Ảnh mờ và không phát hiện được khuôn mặt: yêu cầu chụp lại.
  \item Có thể tắt blur gate; khi đó vẫn cho phép gọi FAS.
  \item Ảnh rõ nét: cho phép gọi FAS.
  \item Quyết định trạng thái: \texttt{RETRY} khi ảnh mờ, \texttt{REJECTED} khi lỗi giấy tờ/tamper,
        \texttt{VERIFIED} khi hợp lệ.
  \item Luồng verify với ảnh selfie mờ phải trả về \texttt{RETRY} kèm \texttt{retake\_reason}.
\end{itemize}

\subsection{Active Liveness}
\textit{Tệp kiểm thử:} \path{tests/test_active_liveness.py}

\textbf{Mục tiêu.} Kiểm tra các trường hợp không có frame, phát hiện chuỗi frame tĩnh (frozen) và xử lý chuỗi
frame có thay đổi.

\begin{itemize}
  \item Danh sách frame rỗng được xử lý đúng.
  \item Các frame giống hệt nhau bị xác định là tĩnh; các frame thay đổi và frame đơn lẻ thì không.
\end{itemize}

\subsection{Tamper Detection}
\textit{Tệp kiểm thử:} \path{tests/test_tamper.py}

\textbf{Mục tiêu.} Kiểm tra kiến trúc TamperNet, hàm loss, letterbox, trích vùng nghi ngờ từ mask, các tín hiệu
heuristic và hợp đồng của engine chỉ dùng heuristic.

\begin{itemize}
  \item Kích thước đầu ra của TamperNet: \texttt{(batch,)} cho logit và \texttt{(batch, 1, 64, 64)} cho mask;
        kích thước ảnh phải chia hết cho kích thước block.
  \item Mô hình có dưới $3\times 10^{6}$ tham số; giá trị loss hữu hạn.
  \item Letterbox/unletterbox giữ box nằm trong phạm vi ảnh gốc.
  \item Các tín hiệu pháp y (forensic): phát hiện lưới JPEG (blockiness), tỷ lệ tần số cao thấp trên ảnh phẳng,
        có/không có box ảnh chân dung; trọng số điểm được chuẩn hóa lại.
  \item Engine chỉ dùng heuristic \textbf{không bao giờ} trả về \texttt{TAMPERED}, chỉ trả về \texttt{SUSPICIOUS}.
  \item Chế độ tắt kiểm tra trả về \texttt{DISABLED}; khi chưa nạp model, engine vẫn trả về đúng hợp đồng.
\end{itemize}

% =====================================================================
\section{Dữ liệu đánh giá, nguồn trích dẫn và giao thức}\label{sec:datasets}
% =====================================================================
Ba bộ dữ liệu công khai được dùng để đánh giá định lượng các mô-đun AI. Nguyên tắc chung:
(i) chỉ dùng dữ liệu để \emph{đánh giá}, không huấn luyện lại trên đó;
(ii) ngưỡng quyết định được chọn trên tập validation (nếu có) và chỉ chấm tập test một lần;
(iii) mọi kết quả đều tái lập được từ các tệp JSON/CSV trong \path{ai/reports/}.
Bảng~\ref{tab:dataset-sources} tóm tắt nguồn gốc và quy mô; Bảng~\ref{tab:dataset-access} ghi
điều kiện truy cập và vị trí lưu trữ cục bộ.

\begin{table}[H]
  \centering
  \caption{Nguồn gốc và quy mô các bộ dữ liệu đánh giá}
  \label{tab:dataset-sources}
  \footnotesize
  \begin{tabular}{@{}p{2.15cm}p{2.0cm}p{5.35cm}p{4.8cm}@{}}
    \toprule
    \textbf{Dataset} & \textbf{Mô-đun} & \textbf{Nguồn / trích dẫn} & \textbf{Quy mô sử dụng}\\
    \midrule
    CASIA-FASD & Passive liveness & Zhang et al. \cite{zhang2012casia}; CBSR (Viện Hàn lâm Khoa học Trung Quốc); gói crop 73 MB & 2.660 frame sau khử trùng lặp (674 thật / 1.986 giả); 30 người / 360 video; $256\times256$\\
    LFW View 1 & So khớp khuôn mặt & Huang et al. \cite{huang2008lfw}; Đại học Massachusetts Amherst & 13.233 ảnh / 5.749 người; dùng 4.992 ảnh duy nhất; 2.200 cặp val + 1.000 cặp test\\
    FantasyID & Phát hiện giả mạo giấy tờ & Korshunov et al. \cite{korshunov2025fantasyid}; Zenodo record 17063366 (v1, 2,6 GB) & Train 1.899 (633 thật / 1.266 giả); val 284 (116 / 168); test 1.385 (300 / 1.085); 911 mask GT\\
    \bottomrule
  \end{tabular}
\end{table}

\begin{table}[H]
  \centering
  \caption{Điều kiện truy cập, giấy phép và vị trí dữ liệu cục bộ}
  \label{tab:dataset-access}
  \footnotesize
  \begin{tabular}{@{}p{2.1cm}p{4.6cm}p{3.4cm}p{4.2cm}@{}}
    \toprule
    \textbf{Dataset} & \textbf{Điều kiện truy cập / giấy phép} & \textbf{Vị trí cục bộ} & \textbf{Liên kết}\\
    \midrule
    CASIA-FASD & Học thuật, phi thương mại; gói crop là bản cắt lại không chính thức & \path{ai/data/archive (1)} & \href{http://www.cbsr.ia.ac.cn/english/FaceAntiSpoofDatabases.asp}{cbsr.ia.ac.cn}\\
    LFW & Mở cho nghiên cứu học thuật & \path{ai/data/lfw-deepfunneled} & \href{http://vis-www.cs.umass.edu/lfw/}{vis-www.cs.umass.edu/lfw}\\
    FantasyID & Giấy phép mở theo công bố Zenodo & \path{ai/data/fantasyid} & \href{https://zenodo.org/records/17063366}{zenodo.org/records/17063366}\\
    \bottomrule
  \end{tabular}
\end{table}

\subsection{CASIA-FASD --- passive liveness}
\textbf{Nguồn.} Bộ dữ liệu do nhóm CBSR (Chinese Academy of Sciences) công bố, mô tả trong
Zhang et al. \cite{zhang2012casia}. Bản gốc gồm 50 chủ thể, 600 video (3 video thật + 9 video tấn công
mỗi chủ thể), ba loại tấn công (ảnh in cong, ảnh in khoét mắt, replay video) và ba mức chất lượng.
Báo cáo này dùng \emph{gói crop khuôn mặt} 73 MB (4.063 ảnh JPEG $256\times256$) --- là bản cắt lại
không chính thức, không có video gốc.

\textbf{Cấu trúc và xử lý.} Tên tệp có dạng
\path{<subject>_<HR?>_<video>.avi_<frame>_<real|fake>.jpg}; nhãn chỉ có \texttt{real}/\texttt{fake},
không tách \textit{print}/\textit{replay}. Gói crop có split không dùng được: 20/20 chủ thể của
\texttt{train\_img} nằm trong \texttt{test\_img} và 1.403 frame trùng khóa
\texttt{(subject, quality, video, frame)}. Vì không huấn luyện trên bộ này, hai split được gộp và
khử trùng lặp, còn \textbf{2.660 frame duy nhất} (674 thật / 1.986 giả), 30 người / 360 video.

\textbf{Giao thức.} Bộ phân loại là MiniFASNetV2 \cite{minivision2020silent}, checkpoint tiền huấn luyện
trên CelebA-Spoof \cite{zhang2020celeba}. Chấm trực tiếp trên crop (\texttt{--no-detect}) để không trộn lỗi bộ phát hiện;
ngưỡng giữ nguyên 0,5 như cấu hình triển khai (\texttt{AI\_FAS\_REAL\_THRESHOLD}), không chọn lại
trên dữ liệu này. Đơn vị đánh giá chính là \textbf{video} (trung bình \texttt{real\_prob} của các frame)
vì các frame trong cùng clip tương quan mạnh; kết quả mức frame được báo cáo kèm. Chỉ số theo
ISO/IEC 30107-3 \cite{iso30107}: APCER, BPCER, ACER và EER.

\subsection{LFW --- so khớp khuôn mặt}
\textbf{Nguồn.} Labeled Faces in the Wild (LFW), phiên bản View 1 deep-funneled, do Huang et al.
\cite{huang2008lfw} công bố; ảnh đã căn mặt và được chuẩn hóa kích thước. Bộ gốc có 13.233 ảnh của
5.749 người.

\textbf{Giao thức.} Dùng đúng cặp chuẩn: 1.100 cặp genuine + 1.100 cặp impostor của DevTrain làm
validation, và 500 + 500 cặp của DevTest làm test. Mỗi ảnh được trích embedding bằng đúng pipeline
triển khai: UltraFace \cite{linzaer2019ultraface} $\rightarrow$ MediaPipe 5 điểm
\cite{lugaresi2019mediapipe} $\rightarrow$ ArcFace $112\times112$ \cite{deng2019arcface}. Điểm của
một cặp là cosine similarity. Ngưỡng EER được chọn trên validation và áp nguyên sang test; ngưỡng
production 0,30 (\texttt{AI\_MATCH\_SCORE}) được báo cáo song song. Cặp có ít nhất một ảnh không phát
hiện được mặt nhận điểm 0 và vẫn được tính vào FAR/FRR chính; báo cáo thêm \texttt{detected\_only}
để tách lỗi phát hiện khỏi chất lượng embedding.

\subsection{FantasyID --- phát hiện giả mạo giấy tờ}
\textbf{Nguồn.} FantasyID do Korshunov et al. \cite{korshunov2025fantasyid} công bố (IJCB 2025),
gồm 13 mẫu thẻ giả tưởng, 13 ngôn ngữ, ảnh bona fide và các tấn công số (face swap, inpaint văn bản,
digital manipulation), kèm vùng bị sửa trong JSON. Bản dùng trong báo cáo là v1 trên Zenodo
(2,6 GB, giấy phép mở).

\textbf{Cấu trúc.} Sau khi chuyển sang layout \path{bona_fide/ forged/ masks/}:
train 1.899 ảnh (633 bona fide / 1.266 forged; 1.091 mask GT), val 284 ảnh (116 / 168), test 1.385 ảnh
(300 / 1.085; 911 mask GT). Tập test có 150 tấn công \texttt{face} và 935 tấn công \texttt{text};
các thiết bị gồm \texttt{huawei}, \texttt{iphone15}, \texttt{scan} và \texttt{iphone15pro}.

\textbf{Giao thức.} Dùng đúng split train/test chính thức. Val là 20\% của train, cắt một lần bằng
\path{torch.Generator().manual_seed(42)} và đóng băng thành \path{models/fantasyid_val.txt}.
Calibration chỉ thấy val để chọn ngưỡng và trọng số fusion; test chỉ được chấm một lần. Báo cáo ba
ngưỡng: 0,55 (triển khai), 0,57 (best-F1 trên val) và 0,675 (FPR 10\% trên val). Chỉ số gồm
F1/Precision/Recall/FPR/FNR tại ngưỡng cố định, AUC, Average Precision (AP), recall theo loại tấn công,
IoU định vị vùng sửa và độ trễ.

\subsection{Các bộ dữ liệu cân nhắc nhưng chưa đưa vào}
Để bảo đảm phạm vi và tính tái lập, báo cáo chưa dùng các bộ sau: Replay-Attack (Idiap) và MSU-MFSD
cho liveness; OULU-NPU và NUAA cho tấn công in/replay; MIDV-2020/MIDV-500 cho OCR giấy tờ;
DocTamper và SIDTD cho tamper. Lý do là cần ký thỏa thuận truy cập, dung lượng lớn, hoặc không có
CCCD Việt Nam. Các bộ này được đề xuất cho giai đoạn tiếp theo (Mục~\ref{sec:future}).

% =====================================================================
\section{Kết quả đánh giá định lượng chi tiết}\label{sec:quant}
% =====================================================================
Bảng~\ref{tab:quant-summary} tóm tắt kết quả chính của ba bộ dữ liệu; các mục tiếp theo trình bày
chi tiết từng bảng, từng ngưỡng và độ trễ/thông lượng (ms, FPS).

\begin{table}[H]
  \centering
  \caption{Tóm tắt kết quả định lượng theo bộ dữ liệu}
  \label{tab:quant-summary}
  \small
  \begin{tabular}{@{}p{2.4cm}p{2.5cm}p{4.2cm}p{4.2cm}@{}}
    \toprule
    \textbf{Dataset} & \textbf{Mô-đun} & \textbf{Chỉ số chính} & \textbf{Kết quả}\\
    \midrule
    CASIA-FASD & Passive liveness & Video EER / APCER / BPCER & 3,33\% / 0,00\% / 10,00\%\\
    LFW & So khớp khuôn mặt & Test EER / FAR@0,30 / Accuracy & 2,00\% / 0,00\% / 98,30\%\\
    FantasyID & Phát hiện giả mạo & Test AUC / AP / F1@0,55 & 0,3417 / 0,7429 / 0,4481\\
    \bottomrule
  \end{tabular}
\end{table}
\subsection{Passive Liveness trên CASIA-FASD}
\subsubsection{Cấu trúc dữ liệu}
\begin{table}[H]
  \centering
  \caption{Cấu trúc dữ liệu CASIA-FASD sau khi gộp và khử trùng lặp}
  \label{tab:casia-composition}
  \small
  \begin{tabular}{@{}ll@{}}
    \toprule
    \textbf{Thành phần} & \textbf{Giá trị}\\
    \midrule
    Số ảnh gốc trong gói crop & 4.063 ảnh JPEG $256\times256$\\
    Sau khử trùng lặp & 2.660 frame duy nhất\\
    Nhãn & 674 thật / 1.986 tấn công\\
    Số người / số video & 30 / 360\\
    Chất lượng video & HR: 30 thật / 90 tấn công; none (LR): 60 / 180\\
    Nhãn tấn công & Chỉ \texttt{real}/\texttt{fake}; không tách print/replay\\
    Ngưỡng quyết định & 0,5 (cố định, không chọn lại)\\
    Đơn vị đánh giá chính & Video (trung bình \texttt{real\_prob} các frame)\\
    \bottomrule
  \end{tabular}
\end{table}

\subsubsection{Kết quả mức frame và mức video}
\begin{table}[H]
  \centering
  \caption{Kết quả passive liveness (MiniFASNetV2) trên CASIA-FASD, TTA bật, ngưỡng 0,5}
  \label{tab:fas-casia}
  \footnotesize
  \resizebox{\textwidth}{!}{%
  \begin{tabular}{@{}lcccccccc@{}}
    \toprule
    \textbf{Mức} & \textbf{N thật} & \textbf{N tấn công} & \textbf{APCER} & \textbf{BPCER} & \textbf{ACER} & \textbf{EER (ngưỡng)} & \textbf{AUC} & \textbf{Median \texttt{real\_prob} (thật / giả)}\\
    \midrule
    Frame & 674 & 1.986 & 0,0136 & 0,1231 & 0,0684 & 0,0639 (0,176) & 0,9874 & 0,9778 / 0,0011\\
    Video & 90 & 270 & 0,0000 & 0,1000 & 0,0500 & 0,0333 (0,264) & 0,9972 & 0,9475 / 0,0039\\
    \bottomrule
  \end{tabular}}
\end{table}

\subsubsection{Điểm vận hành theo mục tiêu BPCER}
\begin{table}[H]
  \centering
  \caption{Điểm vận hành mức frame khi cố định BPCER (TTA bật)}
  \label{tab:fas-operating}
  \small
  \begin{tabular}{@{}lcccc@{}}
    \toprule
    \textbf{Mục tiêu} & \textbf{Ngưỡng} & \textbf{APCER} & \textbf{BPCER} & \textbf{ACER}\\
    \midrule
    BPCER = 5\% & 0,122 & 0,0856 & 0,0504 & 0,0680\\
    BPCER = 10\% & 0,360 & 0,0252 & 0,1009 & 0,0630\\
    \bottomrule
  \end{tabular}
\end{table}

\subsubsection{Kết quả theo chất lượng video}
\begin{table}[H]
  \centering
  \caption{Kết quả CASIA-FASD theo nhóm chất lượng (TTA bật, ngưỡng 0,5)}
  \label{tab:fas-quality}
  \small
  \begin{tabular}{@{}lcccccc@{}}
    \toprule
    \textbf{Nhóm} & \textbf{N thật} & \textbf{N tấn công} & \textbf{APCER} & \textbf{BPCER} & \textbf{ACER} & \textbf{EER (ngưỡng)}\\
    \midrule
    HR & 30 & 90 & 0,0000 & 0,2333 & 0,1167 & 0,0389 (0,179)\\
    none (LR) & 60 & 180 & 0,0000 & 0,0333 & 0,0167 & 0,0167 (0,401)\\
    \bottomrule
  \end{tabular}
\end{table}
Nhóm HR có 7/30 video thật bị từ chối; cỡ mẫu nhỏ nên khoảng tin cậy rộng, chưa đủ để kết luận
``HR khó hơn LR''.

\subsubsection{Ablation TTA và thông lượng}
\begin{table}[H]
  \centering
  \caption{Ảnh hưởng của TTA trên CASIA-FASD (chấm trực tiếp trên crop $256\times256$)}
  \label{tab:fas-tta}
  \footnotesize
  \resizebox{\textwidth}{!}{%
  \begin{tabular}{@{}lcccccc@{}}
    \toprule
    \textbf{Cấu hình} & \textbf{Frame EER} & \textbf{Video BPCER} & \textbf{Video EER (ngưỡng)} & \textbf{Video AUC} & \textbf{Median / P95 (ms)} & \textbf{FPS}\\
    \midrule
    TTA bật (mặc định) & 0,0639 & 0,1000 & 0,0333 (0,264) & 0,9972 & 16,96 / 17,59 & 58,96\\
    TTA tắt & 0,0669 & 0,0778 & 0,0222 (0,303) & 0,9981 & 8,48 / 8,84 & 117,92\\
    \bottomrule
  \end{tabular}}
\end{table}
FPS ở đây là thông lượng bộ phân loại trên crop $256\times256$ (FPS $= 1000 /$ median), không gồm
bước phát hiện khuôn mặt. TTA không mang lại lợi ích rõ ràng: frame EER tốt hơn 0,3 điểm phần trăm
nhưng video EER lại xấu hơn 1,1 điểm phần trăm và độ trễ gấp đôi.

\subsubsection{Độ bền trước suy giảm chất lượng ảnh}
Giữ nguyên checkpoint và ngưỡng 0,5, TTA tắt, áp từng điều kiện lên cả 2.660 crop rồi chấm lại ---
không huấn luyện lại, không chọn ngưỡng trên tập này.

\begin{table}[H]
  \centering
  \caption{Kết quả CASIA-FASD theo 9 điều kiện suy giảm ảnh (mức video, TTA tắt)}
  \label{tab:fas-robustness}
  \footnotesize
  \resizebox{\textwidth}{!}{%
  \begin{tabular}{@{}lcccccc@{}}
    \toprule
    \textbf{Điều kiện} & \textbf{APCER} & \textbf{BPCER} & \textbf{ACER} & \textbf{EER (ngưỡng)} & \textbf{AUC} & \textbf{$\Delta$EER so với clean}\\
    \midrule
    clean & 0,0000 & 0,0778 & 0,0389 & 0,0222 (0,303) & 0,9981 & ---\\
    JPEG q70 & 0,0000 & 0,1000 & 0,0500 & 0,0333 (0,237) & 0,9977 & +0,0111\\
    JPEG q50 & 0,0074 & 0,1111 & 0,0593 & 0,0333 (0,233) & 0,9977 & +0,0111\\
    JPEG q30 & 0,0074 & 0,1111 & 0,0593 & 0,0444 (0,195) & 0,9968 & +0,0222\\
    Gaussian blur $\sigma=1{,}5$ & 0,0000 & 0,2333 & 0,1167 & 0,0778 (0,084) & 0,9780 & +0,0556\\
    Gaussian blur $\sigma=3{,}0$ & 0,0370 & 0,6222 & 0,3296 & 0,2889 (0,034) & 0,8271 & +0,2667\\
    Hạ--nâng 128 px & 0,0000 & 0,1333 & 0,0667 & 0,0648 (0,122) & 0,9900 & +0,0426\\
    Hạ--nâng 64 px & 0,0037 & 0,3556 & 0,1796 & 0,1222 (0,066) & 0,9478 & +0,1000\\
    Giảm sáng $0{,}5\times$ & 0,0074 & 0,0778 & 0,0426 & 0,0241 (0,277) & 0,9976 & +0,0019\\
    \bottomrule
  \end{tabular}}
\end{table}

\begin{table}[H]
  \centering
  \caption{Ảnh hưởng của suy giảm ảnh theo nhóm chất lượng (mức video, TTA tắt)}
  \label{tab:fas-robustness-quality}
  \small
  \begin{tabular}{@{}lcccc@{}}
    \toprule
    \textbf{Điều kiện} & \textbf{BPCER none} & \textbf{BPCER HR} & \textbf{APCER none} & \textbf{APCER HR}\\
    \midrule
    clean & 3,33\% & 16,67\% & 0,00\% & 0,00\%\\
    JPEG q30 & 5,00\% & 23,33\% & 1,11\% & 0,00\%\\
    Blur $\sigma=1{,}5$ & 11,67\% & 46,67\% & 0,00\% & 0,00\%\\
    Blur $\sigma=3{,}0$ & 48,33\% & 90,00\% & 5,00\% & 1,11\%\\
    Hạ--nâng 128 & 3,33\% & 33,33\% & 0,00\% & 0,00\%\\
    Hạ--nâng 64 & 15,00\% & 76,67\% & 0,56\% & 0,00\%\\
    Giảm sáng $0{,}5\times$ & 1,67\% & 20,00\% & 0,00\% & 2,22\%\\
    \bottomrule
  \end{tabular}
\end{table}

\subsubsection{Diễn giải}
\begin{itemize}
  \item \textbf{JPEG không phải nút thắt:} q30 vẫn giữ APCER 0,74\%, BPCER 11,11\%, EER 4,44\% ---
        chỉ tăng khoảng 2,2 điểm phần trăm EER so với ảnh sạch.
  \item \textbf{Mờ là failure mode chính:} $\sigma=1{,}5$ làm BPCER 23,33\% và EER 7,78\%;
        $\sigma=3{,}0$ làm BPCER 62,22\%, EER 28,89\%, APCER 3,70\%, AUC giảm còn 0,8271.
        Nguyên nhân là MiniFASNet dựa vào texture tần số cao; làm mờ xóa dấu vân tay này nên mặt thật
        bị xếp nhầm sang print/replay (false spoof), không phải false accept.
  \item \textbf{Hạ phân giải cũng nguy hiểm:} 128 px cho BPCER 13,33\%, EER 6,48\%; 64 px cho BPCER
        35,56\%, EER 12,22\%, dù APCER vẫn thấp (0,37\%).
  \item \textbf{Thiếu sáng $0{,}5\times$ gần như vô hại:} BPCER giữ 7,78\%, APCER 0,74\%, EER tăng nhẹ
        lên 2,41\% --- nên ưu tiên chống mờ/hạ phân giải hơn là tăng sáng.
  \item \textbf{Nhóm HR nhạy hơn:} blur $\sigma=1{,}5$ làm BPCER HR 46,67\% (so với 11,67\% nhóm
        none); hạ 64 px làm 76,67\% (so với 15,00\%). Nhóm HR có biên quyết định mỏng hơn ngay cả khi
        sạch (16,67\% so với 3,33\%), nên suy giảm đẩy qua ngưỡng.
  \item \textbf{Blur gate đã triển khai:} Laplacian variance trên crop mặt, ngưỡng mặc định 25
        (\path{AI_FAS_BLUR_GATE}); trung vị ảnh thật sạch là 133,8; chỉ 1,3\% ảnh thật sạch dưới 25,
        trong khi 99,9\% ảnh blur $\sigma=1{,}5$ và 100\% ảnh blur $\sigma=3{,}0$ nằm dưới ngưỡng này.
\end{itemize}

\subsection{So khớp khuôn mặt trên LFW}
\subsubsection{Cấu trúc dữ liệu và pipeline}
\begin{table}[H]
  \centering
  \caption{Cấu trúc dữ liệu và cấu hình đánh giá LFW}
  \label{tab:lfw-composition}
  \small
  \begin{tabular}{@{}ll@{}}
    \toprule
    \textbf{Thành phần} & \textbf{Giá trị}\\
    \midrule
    Số ảnh / số người & 13.233 / 5.749\\
    Ảnh duy nhất được dùng & 4.992\\
    Cache hit & 1.408\\
    Ảnh không phát hiện được mặt & 2\\
    Validation (DevTrain) & 2.200 cặp = 1.100 genuine + 1.100 impostor\\
    Test (DevTest) & 1.000 cặp = 500 genuine + 500 impostor\\
    Pipeline trích embedding & UltraFace $\rightarrow$ MediaPipe 5 điểm $\rightarrow$ ArcFace $112\times112$\\
    Điểm của một cặp & Cosine similarity\\
    Ngưỡng EER (val) / production & 0,1615 / 0,30\\
    \bottomrule
  \end{tabular}
\end{table}

\subsubsection{Kết quả trên validation và test}
\begin{table}[H]
  \centering
  \caption{Kết quả so khớp khuôn mặt (ArcFace) trên LFW, ngưỡng quyết định 0,30}
  \label{tab:face-lfw}
  \footnotesize
  \resizebox{\textwidth}{!}{%
  \begin{tabular}{@{}lccccccc@{}}
    \toprule
    \textbf{Tập} & \textbf{Số cặp} & \textbf{FAR} & \textbf{FRR} & \textbf{TAR} & \textbf{Accuracy} & \textbf{EER (ngưỡng)} & \textbf{AUC}\\
    \midrule
    Validation (DevTrain) & 2.200 & 0,0000 & 0,0282 (31/1.100) & 0,9718 & 0,9859 & 0,0145 (0,1615) & 0,9933\\
    Test (DevTest) & 1.000 & 0,0000 & 0,0340 (17/500) & 0,9660 & 0,9830 & 0,0200 (0,1600) & 0,9918\\
    \bottomrule
  \end{tabular}}
\end{table}

\subsubsection{Điểm vận hành và TAR@FAR}
\begin{table}[H]
  \centering
  \caption{Kết quả test tại ngưỡng EER chọn trên validation (0,1615)}
  \label{tab:lfw-eer-threshold}
  \small
  \begin{tabular}{@{}lcccc@{}}
    \toprule
    \textbf{Tập} & \textbf{FAR} & \textbf{FRR} & \textbf{TAR} & \textbf{Accuracy}\\
    \midrule
    Test (1.000 cặp) & 0,0160 (8/500) & 0,0200 (10/500) & 0,9800 & 0,9820\\
    \bottomrule
  \end{tabular}
\end{table}

\begin{table}[H]
  \centering
  \caption{TAR tại các mức FAR mục tiêu (ngưỡng chọn trên validation)}
  \label{tab:lfw-tar-far}
  \small
  \begin{tabular}{@{}lcccc@{}}
    \toprule
    \textbf{FAR mục tiêu} & \textbf{Ngưỡng val} & \textbf{FAR test} & \textbf{TAR test}\\
    \midrule
    1\% & 0,1797 & 0,0160 & 0,9800\\
    0,1\% & 0,2298 & 0,0040 & 0,9760\\
    \bottomrule
  \end{tabular}
\end{table}

\subsubsection{Detected-only và phân bố điểm}
\begin{table}[H]
  \centering
  \caption{Kết quả \texttt{detected\_only} (loại các cặp có ảnh không phát hiện được mặt)}
  \label{tab:lfw-detected}
  \small
  \begin{tabular}{@{}lcccccc@{}}
    \toprule
    \textbf{Tập} & \textbf{Số cặp} & \textbf{Genuine} & \textbf{Impostor} & \textbf{EER} & \textbf{AUC}\\
    \midrule
    Validation & 2.197 & 1.098 & 1.099 & 0,0137 & 0,9942\\
    Test & 1.000 & 500 & 500 & 0,0200 & 0,9918\\
    \bottomrule
  \end{tabular}
\end{table}

\begin{table}[H]
  \centering
  \caption{Phân bố điểm cosine của các cặp LFW}
  \label{tab:lfw-score-dist}
  \small
  \begin{tabular}{@{}lccc@{}}
    \toprule
    \textbf{Tập} & \textbf{Median genuine} & \textbf{Median impostor} & \textbf{Impostor cao nhất}\\
    \midrule
    Validation & 0,592 & 0,0076 & 0,2298\\
    Test & 0,606 & 0,0092 & 0,2362\\
    \bottomrule
  \end{tabular}
\end{table}

\subsubsection{Thông lượng}
Toàn bộ 3.200 cặp (val + test) chạy trong 105,14 giây, tương đương khoảng \textbf{30,4 cặp/giây}
hay \textbf{32,9 ms/cặp}. Con số này gồm đọc ảnh, tiền xử lý và cache, không phải độ trễ thuần của
mô hình embedding.

\subsubsection{Diễn giải}
\begin{itemize}
  \item Ở ngưỡng triển khai 0,30: \textbf{không có false accept} trên 500 cặp impostor test;
        17/500 cặp genuine bị từ chối (FRR 3,4\%), accuracy 98,3\%. Đây là điểm vận hành thiên về an toàn.
  \item Hạ ngưỡng về 0,1615 (EER val) đổi 8 false accept lấy giảm FRR còn 2,0\%; accuracy 98,2\% ---
        gần như cùng mức, nên \textbf{0,30 vẫn hợp lý} cho eKYC.
  \item EER 2,0\% và AUC 0,9918 cho thấy embedding ArcFace tách genuine/impostor tốt; khoảng cách
        trung vị khoảng 0,58 là rất lớn.
  \item \textbf{Không suy rộng sang CCCD:} LFW là ảnh đã căn, frontal, không có glare, perspective
        hay print; selfie--CCCD thật sẽ khó hơn. Ngoài ra checkpoint ArcFace tiền huấn luyện có thể
        đã thấy LFW (không kiểm chứng được), nên đây là \textbf{cận trên lạc quan}.
  \item \textbf{Bug đã sửa:} fallback \path{_aligned_embedding} từng lặp qua crop của
        \path{detect_faces()} thay vì box của \path{detect_face_boxes()}, gây \texttt{TypeError}
        khi MediaPipe không có landmark; đã sửa trong \path{face_service.py} và thêm regression test
        \path{test_aligned_embedding_falls_back_to_detected_box}.
\end{itemize}

\subsection{Phát hiện giả mạo giấy tờ trên FantasyID}
\subsubsection{Cấu trúc dữ liệu và split}
\begin{table}[H]
  \centering
  \caption{Cấu trúc dữ liệu FantasyID sau khi chuẩn bị}
  \label{tab:fantasyid-composition}
  \small
  \begin{tabular}{@{}ll@{}}
    \toprule
    \textbf{Thành phần} & \textbf{Giá trị}\\
    \midrule
    Nguồn & Zenodo record 17063366, phiên bản v1, 2,6 GB\\
    Số mẫu thẻ / ngôn ngữ & 13 / 13\\
    Train & 1.899 ảnh = 633 bona fide + 1.266 forged; 1.091 mask GT\\
    Validation & 284 ảnh = 116 bona fide + 168 forged\\
    Test & 1.385 ảnh = 300 bona fide + 1.085 forged; 911 mask GT\\
    Tấn công trong test & \texttt{face}: 150; \texttt{text}: 935\\
    Thiết bị trong test & \texttt{huawei}, \texttt{iphone15}, \texttt{scan}, \texttt{iphone15pro}\\
    Ngưỡng báo cáo & 0,55 (triển khai) / 0,57 (best-F1 val) / 0,675 (FPR 10\% val)\\
    Trọng số fusion & $w=0{,}70$ (model) + 0,30 (heuristic)\\
    \bottomrule
  \end{tabular}
\end{table}

\subsubsection{Kết quả theo ngưỡng quyết định}
\begin{table}[H]
  \centering
  \caption{Kết quả TamperNet + heuristic trên tập test FantasyID tại ba ngưỡng}
  \label{tab:tamper-fantasyid}
  \footnotesize
  \resizebox{\textwidth}{!}{%
  \begin{tabular}{@{}lccccccccc@{}}
    \toprule
    \textbf{Ngưỡng} & \textbf{F1} & \textbf{Precision} & \textbf{Recall} & \textbf{FPR} & \textbf{FNR} & \textbf{TP} & \textbf{FP} & \textbf{TN} & \textbf{FN}\\
    \midrule
    0,55 (triển khai) & 0,4481 & 0,6710 & 0,3364 & 0,5967 & 0,6636 & 365 & 179 & 121 & 720\\
    0,57 (best-F1 val) & 0,4301 & 0,6603 & 0,3189 & 0,5933 & 0,6811 & 346 & 178 & 122 & 739\\
    0,675 (FPR 10\% val) & 0,2839 & 0,6278 & 0,1834 & 0,3933 & 0,8166 & 199 & 118 & 182 & 886\\
    \bottomrule
  \end{tabular}}
\end{table}

\subsubsection{Chỉ số không phụ thuộc ngưỡng và theo loại tấn công}
\begin{table}[H]
  \centering
  \caption{Chỉ số không phụ thuộc ngưỡng và độ trễ của mô-đun tamper trên FantasyID test}
  \label{tab:tamper-nonthreshold}
  \small
  \begin{tabular}{@{}ll@{}}
    \toprule
    \textbf{Chỉ số} & \textbf{Giá trị}\\
    \midrule
    AUC (fusion) & 0,3417\\
    AUC (model-only) & 0,3190\\
    Average Precision (AP) & 0,7429\\
    IoU định vị (n = 911 mask GT) & trung bình 0,0 / trung vị 0,0\\
    Độ trễ trung vị / P95 & 252,35 / 460,97 ms\\
    Độ trễ trung bình / lớn nhất & 269,95 / 2.263,03 ms\\
    Thông lượng (theo trung vị) & 3,96 ảnh/giây\\
    \bottomrule
  \end{tabular}
\end{table}

\begin{table}[H]
  \centering
  \caption{Recall theo loại tấn công trên FantasyID test, ngưỡng 0,55}
  \label{tab:tamper-attack}
  \small
  \begin{tabular}{@{}lccc@{}}
    \toprule
    \textbf{Loại tấn công} & \textbf{Số mẫu} & \textbf{Recall} & \textbf{Precision so với bona fide}\\
    \midrule
    \texttt{face} (face swap) & 150 & 0,7933 & 1,0000\\
    \texttt{text} (inpaint văn bản) & 935 & 0,2631 & 1,0000\\
    \bottomrule
  \end{tabular}
\end{table}

\subsubsection{Điểm theo thiết bị}
\begin{table}[H]
  \centering
  \caption{Điểm tamper trung bình theo thiết bị trên FantasyID test, ngưỡng 0,55}
  \label{tab:tamper-device}
  \small
  \begin{tabular}{@{}llcccc@{}}
    \toprule
    \textbf{Thiết bị} & \textbf{Nhãn} & \textbf{n} & \textbf{Tamper score TB} & \textbf{Model prob TB} & \textbf{\% dương tính}\\
    \midrule
    \texttt{huawei} & bona fide & 100 & 0,588 & 0,821 & 80,0\%\\
    \texttt{huawei} & forged & 361 & 0,405 & 0,553 & 45,7\%\\
    \texttt{iphone15} & bona fide & 100 & 0,444 & 0,627 & 56,0\%\\
    \texttt{iphone15} & forged & 100 & 0,482 & 0,655 & 62,0\%\\
    \texttt{scan} & bona fide & 100 & 0,374 & 0,527 & 43,0\%\\
    \texttt{scan} & forged & 362 & 0,225 & 0,308 & 24,9\%\\
    \texttt{iphone15pro} & forged & 262 & 0,173 & 0,242 & 18,3\%\\
    \bottomrule
  \end{tabular}
\end{table}
Khoảng cách \textbf{giữa các thiết bị} lớn hơn khoảng cách \textbf{giữa hai lớp} --- bằng chứng định
lượng cho hiện tượng shortcut theo thiết bị.

\subsubsection{Ablation trọng số fusion}
Với $\text{score}(w) = w\cdot\text{model} + (1-w)\cdot\text{signal}$, ngưỡng được chọn lại trên val
cho từng $w$; kết quả test như sau:

\begin{table}[H]
  \centering
  \caption{Ablation trọng số fusion trên FantasyID}
  \label{tab:tamper-ablation}
  \footnotesize
  \resizebox{\textwidth}{!}{%
  \begin{tabular}{@{}lcccccc@{}}
    \toprule
    \textbf{$w_{\text{model}}$} & \textbf{Val AUC} & \textbf{Val F1 @ ngưỡng} & \textbf{Test AUC} & \textbf{Test AP} & \textbf{Test F1 @0,55} & \textbf{FPR @ ngưỡng}\\
    \midrule
    0,0 (heuristic-only) & 0,9603 & 0,8861 @ 0,050 & 0,5172 & 0,8338 & 0,0000 & 0,1433\\
    0,3 & 0,9111 & 0,8713 @ 0,310 & 0,3619 & 0,7624 & 0,0000 & 0,2800\\
    0,5 & 0,8822 & 0,8247 @ 0,495 & 0,3521 & 0,7541 & 0,1081 & 0,3500\\
    0,7 (triển khai) & 0,8619 & 0,8146 @ 0,565 & 0,3417 & 0,7430 & 0,4481 & 0,5967\\
    1,0 (model-only) & 0,8112 & 0,8057 @ 0,815 & 0,3190 & 0,6989 & 0,5070 & 0,5933\\
    \bottomrule
  \end{tabular}}
\end{table}

\subsubsection{Diễn giải}
\begin{itemize}
  \item \textbf{Val sụp khi sang test:} nguyên nhân là domain shift và shortcut thiết bị. Train chỉ có
        \texttt{face\_text} và 3 thiết bị; test là OOD với \texttt{digital\_3},
        \texttt{facedancer} và \path{textdiffuserft_bfei}; bona fide test toàn bộ là
        \texttt{iphone15} --- thiết bị không có trong train; cộng thêm tấn công văn bản khó nhất.
  \item \textbf{AUC < 0,5} là hệ quả trực tiếp của đảo điểm theo thiết bị: trong từng thiết bị hai lớp,
        bona fide có điểm trung bình cao hơn forged, nên xếp hạng toàn cục bị đảo.
  \item \textbf{Heuristic học ``dấu vân tay thiết bị'':} AUC 0,9603 trên val nhưng về ngẫu nhiên
        (0,5172) trên test OOD --- không nên dùng làm tín hiệu quyết định.
  \item \textbf{Model chuyển miền tốt hơn:} AUC/AP cao hơn ở mọi $w<1$; nhưng tại ngưỡng triển khai
        0,55, model-only cho F1 cao hơn (0,5070 so với 0,4481). Nếu buộc phải chọn một cấu hình vận hành,
        $w=1{,}0$ + ngưỡng 0,55 tốt hơn; giữ fusion 0,7 chỉ để giữ khả năng giải thích.
  \item \textbf{Heuristic chỉ nên là thông tin phụ (advisory)}, không tham gia chấm điểm quyết định ---
        đúng với khuyến nghị chuyển tamper sang advisory.
  \item \textbf{IoU = 0} do mask head \emph{không dự đoán} vùng nào, không phải do chất lượng định vị;
        cần phân tích riêng head/regression loss trước khi kết luận.
  \item \textbf{Không được auto-reject:} FPR test 0,5967 nghĩa là gần 60\% người thật bị từ chối nếu
        dùng ngưỡng 0,55 làm quyết định cứng.
\end{itemize}

\subsection{Độ trễ và thông lượng toàn pipeline}
\subsubsection{Độ trễ từng mô-đun}
\begin{table}[H]
  \centering
  \caption{Độ trễ và thông lượng từng mô-đun trên 200 ảnh FantasyID test (đo một luồng)}
  \label{tab:latency}
  \footnotesize
  \resizebox{\textwidth}{!}{%
  \begin{tabular}{@{}lrrrrrl@{}}
    \toprule
    \textbf{Mô-đun} & \textbf{n} & \textbf{Mean (ms)} & \textbf{Median (ms)} & \textbf{P95 (ms)} & \textbf{Max (ms)} & \textbf{FPS (theo median)}\\
    \midrule
    I/O (đọc + RGB) & 400 & 18,96 & 12,85 & 51,43 & 76,13 & 77,82\\
    Face detection (UltraFace) & 200 & 46,89 & 33,01 & 133,55 & 258,41 & 30,29\\
    Landmark (MediaPipe) & 200 & 17,99 & 15,76 & 35,34 & 73,29 & 63,45\\
    Embedding (ArcFace) & 110 & 22,95 & 6,77 & 67,18 & 130,25 & 147,71\\
    Passive Liveness (MiniFASNetV2 + TTA) & 113 & 175,19 & 185,27 & 241,81 & 352,81 & 5,40\\
    OCR (RapidOCR) & 200 & 2.069,50 & 2.329,84 & 3.381,39 & 4.329,14 & 0,43\\
    QR (ảnh mẫu, lặp 20 lần) & 20 & 55,12 & 56,52 & 63,09 & 64,25 & 17,69\\
    Tamper (TamperNet) & 1.385 & 269,95 & 252,35 & 460,97 & 2.263,03 & 3,96\\
    \midrule
    \textbf{Tổng tuần tự (ước tính)} & -- & -- & \textbf{2.892,40} & -- & -- & \textbf{0,35}\\
    \bottomrule
  \end{tabular}}
\end{table}
Ghi chú: FPS $= 1000 /$ median, đo trên một luồng, không batch; đây là thông lượng tham chiếu trên
laptop (RTX 4060 Laptop + ONNX Runtime CPU cho UltraFace/ArcFace/RapidOCR), không phải cam kết SLA
trên server. Tổng tuần tự 2.892,40 ms tương đương khoảng \textbf{2,89 giây/ảnh} và chưa gồm thời gian
upload mạng.

\subsubsection{Độ trễ bộ phân loại FAS trên crop}
\begin{table}[H]
  \centering
  \caption{Độ trễ riêng bộ phân loại MiniFASNetV2 trên crop $256\times256$ (CASIA-FASD)}
  \label{tab:fas-latency}
  \small
  \begin{tabular}{@{}lcccc@{}}
    \toprule
    \textbf{Cấu hình} & \textbf{n} & \textbf{Median (ms)} & \textbf{P95 (ms)} & \textbf{FPS}\\
    \midrule
    TTA bật & 2.660 & 16,96 & 17,59 & 58,96\\
    TTA tắt & 2.660 & 8,48 & 8,84 & 117,92\\
    \bottomrule
  \end{tabular}
\end{table}

\subsubsection{Độ phủ phát hiện trên ảnh chân dung in trên thẻ}
\begin{table}[H]
  \centering
  \caption{Độ phủ phát hiện khuôn mặt và OCR trên 200 ảnh FantasyID test}
  \label{tab:coverage-detect}
  \small
  \begin{tabular}{@{}lccc@{}}
    \toprule
    \textbf{Chỉ số} & \textbf{Số lượng} & \textbf{Tỷ lệ}\\
    \midrule
    UltraFace phát hiện được box & 113/200 & 56,5\%\\
    MediaPipe phát hiện đủ 5 landmark & 110/200 & 55,0\%\\
    Không phát hiện được mặt ở độ phân giải gốc & 87/200 & 43,5\%\\
    OCR thành công & 200/200 & 100,0\%\\
    QR giải mã ảnh mẫu & 20/20 lần lặp & 100,0\%\\
    \bottomrule
  \end{tabular}
\end{table}

\subsubsection{Nhận xét về hiệu năng}
\begin{itemize}
  \item \textbf{OCR là điểm nghẽn chính:} RapidOCR \cite{rapidocr,du2020ppocr} chiếm khoảng \textbf{80,6\%} tổng thời gian tuần tự
        (2.329,84 ms trên 2.892,40 ms). Hạ cạnh dài về khoảng 1.600 px trước khi OCR, crop vùng thẻ hoặc
        chạy OCR song song với các mô-đun khác có thể đưa tổng xuống gần 1 giây.
  \item \textbf{FAS:} 185,27 ms ở mức pipeline (có detect + TTA); trên crop $256\times256$ chỉ còn
        17,0 ms (TTA) hoặc 8,5 ms (không TTA). Tắt TTA giảm khoảng 50\% độ trễ và không làm xấu EER cấp video.
  \item \textbf{Tamper:} trung vị 252,35 ms, P95 460,97 ms, lớn nhất 2.263,03 ms --- có đuôi dài, cần
        đặt timeout và không chặn luồng chính.
  \item \textbf{Độ phủ detector thấp trên mặt in:} 43,5\% ảnh không phát hiện được mặt ở độ phân giải gốc.
        Đây là lỗ hổng thật của nhánh ``mặt trên thẻ''; cần tiling trước khi detect (hàm
        \path{_tile_images} đã có nhưng chỉ dùng trong fallback của \path{_aligned_embedding}).
\end{itemize}

% =====================================================================
\section{Đánh giá độ phủ của bộ kiểm thử}
% =====================================================================
\begin{table}[H]
  \centering
  \caption{Độ phủ kiểm thử theo yêu cầu}
  \label{tab:coverage}
  \small
  \begin{tabular}{@{}p{5.2cm}p{4.3cm}p{4.9cm}@{}}
    \toprule
    \textbf{Yêu cầu kiểm thử} & \textbf{Kiểm thử hiện có} & \textbf{Mức độ phủ}\\
    \midrule
    OCR số CCCD / dự phòng CMND & \path{test_ocr_service.py} & Một phần (chỉ logic trích xuất số)\\
    Khuôn mặt: không mặt / nhiều mặt / so khớp & \path{test_face_service.py} & Một phần (logic và benchmark)\\
    Passive liveness: real / spoof / replay & \path{test_liveness.py} + CASIA-FASD & Một phần (chưa phân loại riêng print/replay)\\
    Active liveness: tư thế / frame tĩnh & \path{test_active_liveness.py} & Thấp (mới có không có frame và frame tĩnh)\\
    Tamper: fusion / an toàn & \path{test_tamper.py} & Rất cao\\
    QR: giải mã / đối chiếu & \path{test_qr_service.py} & Cao (có ảnh mẫu)\\
    Quyết định RETRY / REJECTED / VERIFIED & \path{test_liveness_gate.py} & Cao\\
    \bottomrule
  \end{tabular}
\end{table}

% =====================================================================
\section{Nhận xét và hạn chế}
% =====================================================================
\begin{itemize}
  \item \textbf{Kiểm thử logic:} toàn bộ 44 test đạt, cho thấy phần logic quyết định và tiền xử lý hoạt động
        ổn định; regression về khoảng đầu vào $[0,255]$ của MiniFASNetV2 đã được bảo vệ bằng test.
  \item \textbf{Passive liveness:} EER đạt \num{0.0639} ở mức frame và \num{0.0333} ở mức video trên CASIA-FASD
        (APCER mức video bằng 0). Tuy nhiên BPCER mức video là \num{0.10}, nghĩa là khoảng 10\% mẫu thật bị
        từ chối; điều này được giảm nhẹ nhờ cơ chế \texttt{RETRY} cho phép người dùng chụp lại.
  \item \textbf{So khớp khuôn mặt:} ArcFace đạt EER \num{0.0200} và accuracy \num{0.9830} trên tập test LFW
        tại ngưỡng \num{0.30}, với FAR bằng 0.
  \item \textbf{Phát hiện giả mạo giấy tờ:} kết quả trên FantasyID \textbf{chưa đạt yêu cầu}: AUC chỉ
        \num{0.3417} (thấp hơn mức ngẫu nhiên 0,5), F1 \num{0.4481}, FPR \num{0.5967} và FNR \num{0.6636}.
        Điều này cho thấy mô hình hiện chưa phân biệt được giấy tờ thật và giả trên tập này, có thể do khác biệt
        miền dữ liệu (domain shift) giữa dữ liệu huấn luyện và FantasyID. Do đó mô-đun này chỉ nên đóng vai trò
        tín hiệu phụ trợ, không nên làm căn cứ duy nhất để từ chối giấy tờ. Việc engine chỉ dùng heuristic không
        bao giờ tự trả về \texttt{TAMPERED} là phù hợp với thiết kế an toàn này.
  \item \textbf{Độ trễ:} OCR (RapidOCR) chiếm khoảng \num{80.6}\% tổng thời gian tuần tự của pipeline
        (trung vị \num{2329.84} ms trên tổng \num{2892.40} ms), là điểm nghẽn chính cần tối ưu. Độ trễ
        được đo trên laptop (RTX 4060 Laptop + ONNX Runtime CPU), không đại diện cho server triển khai.
  \item \textbf{CASIA-FASD:} gói crop là bản cắt lại không chính thức, split gốc không dùng được và đã
        được gộp + khử trùng lặp; nhãn không tách print/replay; ảnh $256\times256$ đã căn mặt nên bỏ qua
        bước detect. Vì vậy kết quả cao (video EER 3,33\%) có thể lạc quan so với ảnh in CCCD thật.
  \item \textbf{LFW:} ảnh đã căn, frontal, không có glare/perspective/print; checkpoint ArcFace tiền
        huấn luyện có thể đã thấy LFW (không kiểm chứng được) nên FAR/FRR là \textbf{cận trên lạc quan}.
        Đây là đánh giá recognition trên cặp chuẩn, không phải end-to-end selfie--CCCD.
  \item \textbf{FantasyID:} domain shift và shortcut thiết bị là nguyên nhân chính làm AUC $<0{,}5$;
        bona fide test toàn bộ là \texttt{iphone15} --- thiết bị không có trong train. Kết quả không suy
        rộng trực tiếp cho CCCD Việt Nam; mask head không dự đoán vùng nên IoU $=0$ chưa phản ánh chất
        lượng định vị.
  \item \textbf{Độ bền FAS:} dùng suy giảm tổng hợp (JPEG, Gaussian blur, hạ--nâng phân giải, giảm sáng)
        trên crop đã nén, không phải ảnh chụp điện thoại thật; nhóm HR chỉ có 30 video thật nên khoảng
        tin cậy rộng.
  \item \textbf{Giới hạn khái quát hóa:} các benchmark dùng dữ liệu công khai nên chưa thể suy diễn trực tiếp
        cho CCCD/CMND Việt Nam thật; chưa có thí nghiệm E2E và chưa có người thật tham gia.
\end{itemize}

% =====================================================================
\section{Kết luận và hướng phát triển}
% =====================================================================
\subsection{Kết luận}
Thành phần AI của hệ thống eKYC đã được kiểm chứng bằng 44 test tự động, tất cả đều đạt. Bộ test bao phủ
các quy tắc quyết định và tiền xử lý cốt lõi. Các đánh giá định lượng trên ba bộ dữ liệu công khai
(CASIA-FASD \cite{zhang2012casia}, LFW \cite{huang2008lfw}, FantasyID \cite{korshunov2025fantasyid})
cho kết quả tốt đối với passive liveness và so khớp khuôn mặt; kết quả trên FantasyID cho thấy mô-đun
phát hiện giả mạo giấy tờ còn hạn chế. Báo cáo cũng đã ghi rõ nguồn gốc, quy mô, giấy phép, giao thức
chia tập và kết quả độ trễ/thông lượng (ms, FPS) của từng bộ dữ liệu. Những kết quả này là cơ sở cho
chương cài đặt và kiểm thử của khóa luận, đồng thời định hướng cho các thực nghiệm toàn trình (E2E) và
thực nghiệm trên dữ liệu thật trong tương lai.

\subsection{Hướng phát triển}\label{sec:future}
\begin{enumerate}[label=\arabic*.]
  \item Xây dựng kiểm thử E2E cho toàn bộ luồng FE $\rightarrow$ BE $\rightarrow$ AI $\rightarrow$ Cơ sở dữ liệu.
  \item Bổ sung kiểm thử cho backend (NestJS) và frontend (Next.js).
  \item Bổ sung kiểm thử OCR theo các biến thể chất lượng ảnh (mờ, thiếu/thừa sáng, nghiêng).
  \item Bổ sung kiểm thử tư thế cho active liveness (quay trái/phải).
  \item Thu thập bộ dữ liệu CCCD Việt Nam thật (có sự đồng ý và ẩn danh phù hợp) để benchmark nội miền,
        đồng thời cải thiện mô-đun TamperNet.
\end{enumerate}

\newpage
% =====================================================================
\renewcommand{\refname}{Tài liệu tham khảo}
\begin{thebibliography}{99}
\sloppy

\bibitem{zhang2012casia}
Z. Zhang, J. Yan, S. Liu, Z. Lei, D. Yi, and S. Z. Li,
``A face antispoofing database with diverse attacks,''
in \emph{Proc. IAPR International Conference on Biometrics (ICB)}, 2012, pp. 26--31.

\bibitem{huang2008lfw}
G. B. Huang, M. Mattar, T. Berg, and E. Learned-Miller,
``Labeled Faces in the Wild: A database for studying face recognition in unconstrained environments,''
in \emph{Workshop on Faces in Real-Life Images, ECCV}, 2008.

\bibitem{korshunov2025fantasyid}
P. Korshunov, A. Mohammadi, A. Vidit, C. Ecabert, and S. Marcel,
``FantasyID: A dataset for detecting digital manipulations of ID-documents,''
in \emph{Proc. IEEE International Joint Conference on Biometrics (IJCB)}, 2025.
Dataset v1: \url{https://zenodo.org/records/17063366}.

\bibitem{deng2019arcface}
J. Deng, J. Guo, N. Xue, and S. Zafeiriou,
``ArcFace: Additive angular margin loss for deep face recognition,''
in \emph{Proc. IEEE/CVF Conference on Computer Vision and Pattern Recognition (CVPR)}, 2019, pp. 4690--4699.

\bibitem{zhang2020celeba}
Y. Zhang, Z. Yin, Y. Li, G. Yin, J. Yan, J. Ma, and Z. Liu,
``CelebA-Spoof: Large-scale face anti-spoofing dataset with rich annotations,''
in \emph{Proc. European Conference on Computer Vision (ECCV)}, 2020.

\bibitem{minivision2020silent}
minivision-ai,
``Silent-Face-Anti-Spoofing: MiniFASNetV2 implementation and pretrained weights,''
GitHub repository, 2020.
\url{https://github.com/minivision-ai/Silent-Face-Anti-Spoofing}.

\bibitem{linzaer2019ultraface}
Linzaer,
``Ultra-Light-Fast-Generic-Face-Detector-1MB (UltraFace),''
GitHub repository, 2019.
\url{https://github.com/Linzaer/Ultra-Light-Fast-Generic-Face-Detector-1MB}.

\bibitem{lugaresi2019mediapipe}
C. Lugaresi, J. Tang, H. Nash, C. McClanahan, E. Uboweja, M. Hays, F. Zhang, C.-L. Chang, M. G. Yong,
J. Lee, W.-T. Chang, W. Hua, M. Georg, and M. Grundmann,
``MediaPipe: A framework for building perception pipelines,''
\emph{arXiv preprint arXiv:1906.08172}, 2019.

\bibitem{du2020ppocr}
Y. Du, C. Li, R. Guo, X. Yin, W. Liu, J. Zhou, Y. Bai, Z. Yu, Y. Yang, Q. Dang, and H. Wang,
``PP-OCR: A practical ultra lightweight OCR system,''
\emph{arXiv preprint arXiv:2009.09941}, 2020.

\bibitem{rapidocr}
RapidAI,
``RapidOCR: A cross-platform OCR toolkit based on ONNX Runtime,''
GitHub repository.
\url{https://github.com/RapidAI/RapidOCR}.

\bibitem{iso30107}
ISO/IEC 30107-3:2017,
\emph{Information technology --- Biometric presentation attack detection --- Part 3: Testing and reporting},
International Organization for Standardization, 2017.

\end{thebibliography}

% =====================================================================
\appendix
\section{Danh sách đầy đủ các test case (pytest node ID)}
% =====================================================================
{\footnotesize
\begin{longtable}{@{}p{3.6cm}p{6.8cm}p{4.2cm}@{}}
  \caption{Danh sách 44 test case}\label{tab:appendix}\\
  \toprule
  \textbf{Tệp} & \textbf{Tên test} & \textbf{Mô tả}\\
  \midrule
  \endfirsthead
  \toprule
  \textbf{Tệp} & \textbf{Tên test} & \textbf{Mô tả}\\
  \midrule
  \endhead
  \bottomrule
  \endfoot

  \path{test_active_liveness.py} & \path{test_analyze_no_frames} & Danh sách frame rỗng\\
   & \path{test_is_frozen_true_for_identical_frames} & Phát hiện frame tĩnh\\
   & \path{test_is_frozen_false_for_changing_frames} & Frame thay đổi không bị coi là tĩnh\\
   & \path{test_is_frozen_false_for_single_frame} & Frame đơn lẻ không bị coi là tĩnh\\
  \midrule
  \path{test_face_service.py} & \path{test_detector_uses_model_input_size} & Kích thước đầu vào detector\\
   & \path{test_preprocess_shape} & Kích thước tensor đầu vào\\
   & \path{test_postprocess_decodes_and_suppresses} & Giải mã box và NMS\\
   & \path{test_nms_keeps_distinct_boxes} & Giữ các box tách biệt\\
   & \path{test_detect_face_boxes_runs} & Chạy detector thực tế\\
   & \path{test_aligned_embedding_falls_back_to_detected_box} & Dự phòng khi thiếu landmark\\
  \midrule
  \path{test_liveness.py} & \path{test_get_new_box_stays_within_image} & Box mở rộng nằm trong ảnh\\
   & \path{test_crop_face_returns_target_size} & Kích thước ảnh crop\\
   & \path{test_predict_uses_0_255_input_range} & Regression đầu vào MiniFASNet\\
   & \path{test_predict_returns_three_class_probabilities} & Ba xác suất, tổng bằng 1\\
  \midrule
  \path{test_liveness_gate.py} & \path{test_low_blur_returns_retake_without_calling_fas} & Blur gate, không gọi FAS\\
   & \path{test_blurry_frame_without_detectable_face_returns_retake} & Ảnh mờ và không có mặt\\
   & \path{test_quality_gate_can_be_disabled} & Tắt blur gate\\
   & \path{test_sharp_image_reaches_fas} & Ảnh rõ được đưa vào FAS\\
   & \path{test_decide_status_retry_versus_reject} & Quyết định RETRY/REJECTED\\
   & \path{test_verify_returns_retry_for_blurry_selfie} & Luồng verify với selfie mờ\\
  \midrule
  \path{test_ocr_service.py} & \path{test_prefers_twelve_digit_cccd} & Nhận số CCCD 12 chữ số\\
   & \path{test_prefers_twelve_over_nine} & Ưu tiên CCCD hơn CMND\\
   & \path{test_falls_back_to_nine_digit_cmnd} & Dự phòng sang CMND\\
   & \path{test_returns_none_when_absent} & Không có số thì trả null\\
  \midrule
  \path{test_qr_service.py} & \path{test_parse_with_legacy_id} & Parse QR có số CMND cũ\\
   & \path{test_parse_without_legacy_id} & Parse QR không có số CMND cũ\\
   & \path{test_compare_match_ignores_accents_and_separators} & MATCH, bỏ qua dấu và ký tự phân tách\\
   & \path{test_compare_mismatch_document_number} & MISMATCH số giấy tờ\\
   & \path{test_compare_partial_when_only_name_differs} & PARTIAL khi chỉ khác họ tên\\
   & \path{test_compare_not_found_when_no_qr} & NOT\_FOUND khi không có QR\\
   & \path{test_compare_not_found_when_no_ocr} & NOT\_FOUND khi không có OCR\\
   & \path{test_decode_fixture_image} & Giải mã ảnh mẫu\\
  \midrule
  \path{test_tamper.py} & \path{test_tamper_net_shapes} & Kích thước đầu ra TamperNet\\
   & \path{test_tamper_net_requires_size_divisible_by_block} & Kích thước chia hết cho block\\
   & \path{test_tamper_loss_is_finite} & Loss hữu hạn\\
   & \path{test_letterbox_roundtrip_box} & Letterbox/unletterbox\\
   & \path{test_mask_to_regions_finds_component} & Trích vùng từ mask\\
   & \path{test_blockiness_detects_jpeg_grid} & Phát hiện lưới JPEG\\
   & \path{test_low_hf_fraction_high_on_flat_image} & Tỷ lệ tần số cao trên ảnh phẳng\\
   & \path{test_forensic_signals_with_and_without_photo_box} & Tín hiệu pháp y có/không box ảnh\\
   & \path{test_signal_score_renormalises_weights} & Chuẩn hóa lại trọng số\\
   & \path{test_heuristic_engine_never_rejects_alone} & Heuristic không tự từ chối\\
   & \path{test_analyze_without_model_returns_contract} & Hợp đồng khi chưa nạp model\\
   & \path{test_analyze_disabled} & Chế độ DISABLED\\
\end{longtable}}

\end{document}